% 示例阅读顺序：查看实际效果，在源码中复制对应 frame，再按需替换内容。
\section{Block 与注释}
\subsection{三色块与组合}
\begin{frame}{Block 样式示例}
  \begin{block}{普通说明：block}
    用于研究问题、定义或方法说明。正文可以包含\textbf{粗体}与公式，例如 $E=mc^2$。
  \end{block}

  \begin{alertblock}{注意事项：alertblock}
    用于假设、限制和需要关注的结论；可用 \alert{强调文字}标出关键条件。
  \end{alertblock}

  \begin{exampleblock}{具体示例：exampleblock}
    用于算例、实验设置或应用场景。例如：对三组数据采用相同指标进行比较。
  \end{exampleblock}
\end{frame}

% 三色环境：block / alertblock / exampleblock；分栏使用 [T,onlytextwidth]。
\begin{frame}{Block 中的列表与双栏}
  \begin{columns}[T,onlytextwidth]
    \column{0.48\textwidth}
    \begin{block}{方法步骤}
      \begin{enumerate}
        \item 明确研究对象。
        \item 统一评价指标。
        \item 对比实验结果。
      \end{enumerate}
    \end{block}

    \column{0.48\textwidth}
    \begin{alertblock}{比较条件}
      \begin{itemize}
        \item 使用相同数据集。
        \item 记录参数设置。
        \item 报告适用范围。
      \end{itemize}
    \end{alertblock}
  \end{columns}

  \begin{exampleblock}{带公式的示例}
    样本均值为 $\bar{x}=\frac{1}{n}\sum_{i=1}^{n}x_i$。
    对样本 $1,2,3$，有 $\bar{x}=2$。
  \end{exampleblock}
\end{frame}

% 空字符串和纯空格均视为无标题；三类环境保留各自的正文背景色。
\begin{frame}{无标题 Block}
  \begin{block}{}
    普通信息：空标题不产生标题栏，直接显示正文。
  \end{block}
  \begin{alertblock}{   }
    注意条件：纯空格标题也不会产生空白标题栏。
  \end{alertblock}
  \begin{exampleblock}{}
    应用示例：支持 $x>0$、\textbf{粗体}和多行文字。
  \end{exampleblock}
\end{frame}

\subsection{独立段落注释}
% 用法：\inknote{补充说明}；结束后恢复正文样式。
\begin{frame}[fragile]{正文中的段落注释}
  这是正常正文，用于陈述主要观点和分析结果。
  \inknote{这是一段补充注释：可以解释假设、范围或阅读提示。
  它没有标题栏，采用深蓝细竖线、淡灰直角背景与小字号深灰衬线文字，与正文自然区分。}
  注释结束后，正文恢复原来的字号、颜色和字重。
  \begin{verbatim}
\inknote{这里填写注释内容。}
  \end{verbatim}
\end{frame}

\begin{frame}{分栏注释与来源说明}
  \begin{columns}[T,onlytextwidth]
    \column{0.48\textwidth}
      \textbf{方法说明}\par\smallskip
      这里放置正文段落。
      \inknote{注释自动适应当前栏宽，也可包含 $n>0$、\textbf{粗体}和链接。

      空行可分隔注释内部的段落。}
    \column{0.48\textwidth}
      \textbf{结果说明}\par\smallskip
      这里放置另一段正文。
      \inknote{inknote 用于解释正文；inksource 用于简短的图表出处。
      两者的语义和视觉层级不同。}
  \end{columns}
  \inksource{来源：此处填写数据来源或参考文献。}
\end{frame}
